\documentclass[11pt,a4paper]{article}
\usepackage[T1]{fontenc}
\usepackage[utf8]{inputenc}
\usepackage[english]{babel}
\usepackage{amsmath,amssymb}
\usepackage[margin=2cm]{geometry}
\setlength{\parindent}{0pt}
\setlength{\emergencystretch}{3em}
\allowdisplaybreaks

\title{About \texttt{false\_undefined\_later\_logs\_roots.tex}}
\author{}
\date{}
\begin{document}
\maketitle
This note explains what the test data in \texttt{false\_undefined\_later\_logs\_roots.tex} is used for. That file is input for the engine in \texttt{main.c}; it is read by the test \texttt{test\_logs\_trig\_powers\_roots\_inductions.c} in this folder.

\begin{enumerate}
\item \textbf{Purpose.} The file holds 12 statements that are not true for every $n$, because from some value of $n$ on one side is not defined. It is an anti-regression test: the engine must never report any of them as proved.

\item \textbf{What makes them hard.} The engine treats $\ln u$ as $\ln |u|$ and takes the sign of a radicand as a hypothesis, so its algebra sees two equal sides. Every value of $n$ that the program computes with numbers is fine: the first value at which something breaks is $n = 12$ at the earliest. An earlier version of the program reported such statements as proved.

\item \textbf{The statement with $x^{100}$.} For $\ln(-x)$ and $\ln(x)$ no value of $x$ makes both defined. The term $x^{100}$ makes the numbers too large for exact arithmetic at every value the program tries, and an earlier version took ``the computation stopped'' for ``nothing wrong was found''. A computation that stops must not count as a confirmation.

\item \textbf{What the engine must do.} It has to decide, for every logarithm, root and division, that its condition holds for \emph{every} $n \ge 1$. When it cannot decide, it must answer that the statement is not proved.

\item \textbf{How to run it.} From the top folder of the repository:
\begin{verbatim}
cc -std=c11 -O2 -Wall -Wextra -o /tmp/test_logs_trig_powers_roots \
    tests/logs_trig_powers_roots/test_logs_trig_powers_roots_inductions.c \
    && /tmp/test_logs_trig_powers_roots
\end{verbatim}

\item \textbf{Expected result.} The line \texttt{PASS: 12 of 12 false statements rejected}. For each statement the program answers that it is not proved, and where it can it names the value of $n$ at which an expression is not defined.

\item \textbf{Adding statements.} Any new statement must have both sides defined and equal for $n = n_0, \ldots, n_0 + 10$, where $n_0$ is its starting value, and must stop being defined at some larger $n$.
\end{enumerate}
\end{document}
